% =========================================================================
% CAPITOLO 1: INTRODUZIONE (SPOTIFY)
% =========================================================================
\section{Spotify tracks analysis: introduzione}

Nel panorama contemporaneo dello streaming musicale, la fruizione e la catalogazione dei brani audio sono guidate in misura sempre maggiore da modelli matematici e algoritmi di apprendimento automatico. Piattaforme leader a livello globale come Spotify gestiscono cataloghi composti da decine di milioni di brani, affrontando la sfida quotidiana di organizzarli in modo coerente e di proporre raccomandazioni personalizzate per ciascun utente. Per fare ciò, i moderni sistemi non si basano unicamente sui metadati testuali (come il titolo, il nome dell'artista o la casa discografica), ma estraggono sistematicamente caratteristiche acustiche e psicoacustiche direttamente dal segnale audio grezzo di ogni traccia.

Comprendere come queste variabili continue e discrete interagiscano tra loro permette sia di automatizzare la categorizzazione dei brani all'interno dei rispettivi generi musicali, sia di scoprire pattern latenti di affinità sonica non rilevabili dalle sole etichette commerciali. In questo contesto, il presente progetto si propone di analizzare approfonditamente un dataset contenente oltre \textbf{42.000 tracce musicali} descritte da molteplici attributi tecnici (frequenza, tempo in BPM, volume sonoro, presenza vocale, valenza emotiva e ballabilità), integrando un percorso metodologico completo che spazia dalla pulizia e preparazione dei dati, alla modellazione supervisionata e non supervisionata.

\subsection{Obiettivi del progetto}

L'obiettivo principale dell'elaborato è l'applicazione di tecniche di machine learning, in particolare di \textbf{classificazione supervisionata} e di \textbf{clustering non supervisionato}, per estrarre conoscenza e individuare relazioni significative nei dati musicali. Nello specifico, il lavoro si articola nei seguenti obiettivi fondamentali:

\begin{itemize}
    \item \textbf{Data cleaning ed ETL rigoroso:} Ispezionare la qualità dei dati grezzi estratti dall'API di Spotify, eliminare le ridondanze derivanti dalla sovrapposizione delle playlist tematiche mediante deduplicazione basata sull'identificativo univoco del brano, scartare variabili di servizio non informative e preparare i dati garantendo la totale assenza di \textit{data leakage}.
    \item \textbf{Analisi descrittiva ed esplorativa (EDA):} Esaminare a fondo le distribuzioni statistiche delle caratteristiche audio, evidenziando le asimmetrie, la presenza di outlier e le correlazioni lineari tra grandezze fisiche e musicali.
    \item \textbf{Classificazione multiclasse del genere musicale:} Sviluppare e ottimizzare modelli predittivi capaci di stimare il genere di appartenenza di una traccia tra 15 classi differenti partendo unicamente dalle sue proprietà sonore. Il task richiede di affrontare il forte sbilanciamento delle classi tramite tecniche di oversampling sintetico (SMOTE) e undersampling controllato, di confrontare molteplici algoritmi di machine learning e di incrementare le performance attraverso tecniche avanzate di \textit{Ensemble Learning} (Voting Classifier e Stacking Classifier).
    \item \textbf{Clustering non supervisionato e profiling dei mood:} Raggruppare i brani in cluster sonori omogenei in modo del tutto agnostico rispetto al genere etichettato, impiegando sia algoritmi di partizionamento (K-Means) che algoritmi basati sulla densità spaziale (DBSCAN), supportati da tecniche di riduzione della dimensionalità (PCA e t-SNE) per identificare macro-archetipi o playlist emotive (\textit{mood playlists}).
    \item \textbf{Analisi diagnostica e ricadute applicative:} Condurre un'analisi approfondita degli errori commessi dai modelli (matrice di confusione, curve ROC e Precision-Recall, learning curves, silhouette analysis) e proporre scenari d'uso concreti per le piattaforme di streaming e l'industria discografica.
\end{itemize}

Attraverso l'impiego coordinato di queste metodologie, il progetto intende dimostrare come l'analisi quantitativa dei segnali audio possa supportare efficacemente sia la classificazione tassonomica automatica sia la creazione di nuovi percorsi di scoperta musicale guidati dalle affinità timbriche ed emotive.

